\begin{table}[!htbp]\centering
\caption{Variantes de especificación de los tres resultados}\label{tab:short_specs}
\begin{threeparttable}\small
\setlength{\tabcolsep}{4pt}
\begin{tabular}{lrrr}
\toprule
Variante & Salario (log) & Formalidad (pp) & Autoempleo (pp) \\
\midrule
Base & -0.021 (0.012) & -4.642 (0.739) & 3.630 (0.595) \\
Sin controles & -0.008 (0.012) & -4.198 (0.749) & 4.222 (0.590) \\
Excluye 2021T1 & -0.023 (0.012) & -4.196 (0.828) & 3.135 (0.500) \\
Tendencia por departamento & -0.022 (0.012) & -4.043 (0.729) & 3.439 (0.620) \\
Edad de 25 a 55 años & -0.030 (0.015) & -4.192 (0.765) & 2.884 (0.659) \\
Corte educativo alternativo & -0.015 (0.012) & -6.483 (0.633) & 4.998 (0.663) \\
EF ocupación y rama & -0.007 (0.011) & -3.645 (0.583) & 2.715 (0.472) \\
Incluye 2022T2 & -0.018 (0.012) & -4.367 (0.724) & 3.541 (0.617) \\
Excluye agricultura & -0.022 (0.012) & -3.763 (0.726) & 2.794 (0.531) \\
Incluye sector público & -0.002 (0.010) & -3.712 (0.801) & 2.995 (0.577) \\
Tendencia por grupo educativo & -0.021 (0.019) & -0.353 (0.970) & 0.607 (0.894) \\
EF departamento por trimestre & -0.020 (0.013) & -4.132 (0.770) & 3.399 (0.624) \\
\bottomrule
\end{tabular}
\begin{tablenotes}\footnotesize
\item Coeficientes DiD y errores agrupados por departamento entre paréntesis. Las variantes modifican una característica de la especificación base; los tamaños de muestra completos están en los archivos de replicación. El corte alternativo compara secundaria incompleta o menos con superior no universitaria completa o más. Ocupación y rama pueden ser variables posteriores al tratamiento. La tendencia por grupo se estima en toda la ventana y puede absorber tanto una deriva como parte de un cambio persistente; su resultado no valida ninguna de las dos explicaciones por sí solo.
\end{tablenotes}
\end{threeparttable}\end{table}
